\subsection{Prime ideals associated with a graded module}

PROPOSITION 1. Let \(\Delta\) be a torsion-free commutative group, A a graded ring of type \(\Delta\) and M a graded A-module of type \(\Delta\) . Every prime ideal associated with M is graded and is the annihilator of a homogeneous element of M.

We know that \(\Delta\) can be given a total order structure compatible with its group structure (Algebra, Chapter II, § 11, no. 4, Lemma 2). Let p be a prime ideal associated with M, the annihilator of an element \(x \in M\) , and let \((x_{i})_{i \in \Delta}\) be the family of homogeneous components of x; let \(i(1) < i(2) < \cdots < i(r)\) be the values of i for which \(x_{i} \neq 0\) . Consider an element \(a \in p\) and let \((a_{i})_{i \in \Delta}\) be the family of its homogeneous components; we shall prove that \(a_{i} \in p\) for all \(i \in \Delta\) , which will show that p is a graded ideal.

We argue by induction on the number of indices \(i\) such that \(a_{i} \neq 0\) . Our assertion is obvious if this number is 0; if not, let \(m\) be the greatest of the indices \(i\) for which \(a_{i} \neq 0\) ; if we prove that \(a_{m} \in \mathfrak{p}\) , the induction hypothesis applied to \(a - a_{m}\) will give the conclusion. Now, \(ax = 0\) ; for all \(j \in \Delta\) , using the fact that the homogeneous component of degree \(m + j\) of \(ax\) is 0, we obtain \(\sum_{i \in \Delta} a_{m - i} x_{j + i} = 0\) ; we conclude that \(a_{m} x_{j}\) is a linear combination of the \(x_{i}\) of indices \(i > j\) . In particular therefore \(a_{m} x_{i(r)} = 0\) , whence, by descending induction on \(n < r\) , \(a_{m}^{r - n + 1} x_{i(n)} = 0\) . Then \(a_{m}^{r} x = 0\) , whence \(a_{m}^{r} \in \mathfrak{p}\) and, as \(\mathfrak{p}\) is prime, \(a_{m} \in \mathfrak{p}\) .

We now show that \(\mathfrak{p}\) is the annihilator of a homogeneous element of M. Let us write \(\mathfrak{b}_n = \operatorname{Ann}(x_{i(n)})\) for \(1 \leqslant n \leqslant r\) . For every homogeneous element \(b\) of \(\mathfrak{p}\) and all \(n\) the homogeneous component of \(bx\) of degree \(i(n) + \deg(b)\) is \(bx_{i(n)}\) , hence \(bx_{i(n)} = 0\) and therefore \(b \in \mathfrak{b}_n\) ; as \(\mathfrak{p}\) is generated by its homogeneous elements, \(\mathfrak{p} \subset \mathfrak{b}_n\) . On the other hand, clearly \(\bigcap_{n=1}^{r} \mathfrak{b}_n \subset \operatorname{Ann}(x) = \mathfrak{p}\) ; as \(\mathfrak{p}\) is prime, there exists an \(n\) such that \(\mathfrak{b}_n \subset \mathfrak{p}\) (Chapter II, § 1, no. 1, Proposition 1), whence \(\mathfrak{b}_n = \mathfrak{p} = \operatorname{Ann}(x_{i(n)})\) , which completes the proof.

COROLLARY. For every (necessarily graded) prime ideal p associated with a graded A-module M, there exists an index \(k \in \Delta\) such that the graded A-module \((\mathbf{A}/\mathfrak{p})(k)\) obtained from the graded A-module A/p by diminishing the degrees by k (Algebra, Chapter II, § 11, no. 2) is isomorphic to a graded submodule of M.

With the notation of the proof of Proposition 1, consider the homomorphism obtained, by taking quotients, from the homomorphism \(a \mapsto ax_{i(n)}\) of A to M; the latter is a graded homomorphism of degree \(i(n)\) and hence it gives on taking quotients a graded bijective homomorphism of degree \(i(n)\) of \(A / p\) onto a graded submodule of \(M\) .

PROPOSITION 2. Let \(\Delta\) be a torsion-free commutative group, A a graded Noetherian ring of type \(\Delta\) and M a graded finitely generated A-module of type \(\Delta\) . There exists a composition series \((\mathbf{M}_i)_{0 \leqslant i \leqslant n}\) consisting of graded submodules of M such that for \(0 \leqslant i \leqslant n - 1\) the graded module \(\mathbf{M}_i / \mathbf{M}_{i+1}\) is isomorphic to a shifted graded module \((\mathbf{A} / \mathfrak{p}_i)(k_i)\) , where \(\mathfrak{p}_i\) is a graded prime ideal of A and \(k_i \in \Delta\) .

It is sufficient to retrace the argument of § 1, no. 4, Theorem 1 taking on this occasion \(\mathfrak{E}\) to be the set of graded submodules of M with a composition series with the properties of the statement; we conclude using the Corollary to Proposition 1.

